%  LaTeX support: latex@mdpi.com
%  Reformatted from Elsevier elsarticle to MDPI format
%  v2: Added Python scripts, restructured Section 2, enriched Results

%=================================================================
\documentclass[environments,article,submit,pdftex,moreauthors]{Definitions/mdpi}

%=================================================================
% Journal selection -- change to target MDPI journal as needed
% Suitable candidates: soilsystems, ijerph, environments, sustainability, toxics
%=================================================================

% MDPI internal commands
\firstpage{1}
\makeatletter
\setcounter{page}{\@firstpage}
\makeatother
\pubvolume{1}
\issuenum{1}
\articlenumber{0}
\pubyear{2026}
\copyrightyear{2026}
%\externaleditor{Academic Editor: Firstname Lastname}
\datereceived{23 May 2026}
\daterevised{}
\dateaccepted{}
\datepublished{}
\hreflink{https://doi.org/}
%\doinum{}

%=================================================================
% Packages
%=================================================================
\graphicspath{{figures/}}
\usepackage{subcaption}
\usepackage{listings}
\usepackage{xcolor}
\usepackage{gensymb}
\usepackage{tabularx}
\newcolumntype{Y}{>{\centering\arraybackslash}X}
\usepackage{amsmath,mathtools,amsthm}
    \setcounter{MaxMatrixCols}{15}
\usepackage{array}
\usepackage{amssymb}
\usepackage{amsfonts}
\usepackage{float}
\usepackage{chngcntr}
\counterwithin*{equation}{section}
\usepackage{longtable}
\usepackage{multirow}
\usepackage{booktabs}
\usepackage{adjustbox}

%=================================================================
% Python listings style with full syntax colouring
%=================================================================
\definecolor{codebg}{RGB}{245,247,250}
\definecolor{codeframe}{RGB}{180,190,210}
\definecolor{kwcolor}{RGB}{0,95,185}       % keywords  -- deep blue
\definecolor{builtincolor}{RGB}{160,32,240}% built-ins -- purple
\definecolor{strcolor}{RGB}{180,60,0}      % strings   -- brick
\definecolor{cmtcolor}{RGB}{80,130,80}     % comments  -- green
\definecolor{numcolor}{RGB}{0,130,100}     % numbers   -- teal
\definecolor{fncolor}{RGB}{0,100,160}      % functions -- blue

\lstdefinestyle{PythonStyle}{
    language=Python,
    backgroundcolor=\color{codebg},
    frame=single,
    framerule=0.6pt,
    rulecolor=\color{codeframe},
    basicstyle=\ttfamily\footnotesize,
    keywordstyle=\color{kwcolor}\bfseries,
    stringstyle=\color{strcolor},
    commentstyle=\color{cmtcolor}\itshape,
    numberstyle=\tiny\color{gray},
    numbers=left,
    numbersep=6pt,
    showstringspaces=false,
    breaklines=true,
    breakatwhitespace=false,
    tabsize=4,
    captionpos=b,
    keepspaces=true,
    xleftmargin=14pt,
    morekeywords={as,with,yield,lambda,assert,pass,del,
                  import,from,raise,try,except,finally,
                  class,def,return,if,elif,else,for,while,
                  in,not,and,or,is,True,False,None},
    emph={print,range,len,list,dict,set,tuple,int,float,str,
          bool,type,enumerate,zip,map,filter,sorted,sum,min,max,
          open,close,append,extend,fit,predict,transform,
          score,reshape,ravel,flatten},
    emphstyle=\color{fncolor},
}
\lstset{style=PythonStyle}

%=================================================================
\Title{Python-Powered Environmental Intelligence: AI and ML Workflows for Soil Pollution Assessment and SES Modelling}

\TitleCitation{Python-Powered Environmental Intelligence: AI and ML Workflows for Soil Pollution Assessment and SES Modelling}

% Author ORCID
\newcommand{\orcidauthorA}{0000-0002-5759-1089}

% Authors
\Author{Polina Lemenkova $^{1,2,3,4}$*\orcidA{}}

\AuthorNames{Polina Lemenkova}
\AuthorCitation{Lemenkova, P.}

\address{%
$^{1}$ \quad Bureau de Recherches G\'eologiques et Mini\`eres (BRGM), Orl\'eans, France\\
$^{2}$ \quad Le Studium, Institut d'Etudes Avanc\'ees du Val de Loire,
    1 rue Dupanloup 45000, Orl\'eans, France.\\
$^{3}$ \quad Le laboratoire PRISME, Universit\'e d'Orl\'eans, Institut National
    des Sciences Appliqu\'ees de Centre-Val de Loire (INSA CVL, UR 4229),
    Orl\'eans, France.\\
$^{4}$ \quad La Technopole d'Orl\'eans, 1 Avenue du Champ de Mars,
    45100 Orl\'eans, France.
}

\corres{Correspondence: polina.lemenkova@tech-orleans.fr}

%=================================================================
% Abstract
%=================================================================
\abstract{Soil pollution represents one of the most pressing global
environmental challenges, driven by industrialization, intensive agriculture,
urban expansion, mining, and the excessive application of fertilizers,
pesticides, and synthetic chemicals. This integrative review synthesizes
evidence from more than 100 peer-reviewed studies to evaluate computational,
spectroscopic, and artificial intelligence (AI)-driven methodologies for soil
pollution assessment across five major contaminant groups: heavy metals,
pesticides, microplastics, per- and polyfluoroalkyl substances (PFAS), and
excess macronutrients. Advanced spectroscopic approaches---Fourier transform
infrared spectroscopy (FTIR), X-ray photoelectron spectroscopy (XPS),
hyperspectral imaging, and synchrotron radiation techniques---are benchmarked
against data-driven machine learning (ML) frameworks implemented in Python.
Seven annotated Python workflows are presented, demonstrating the application
of Random Forest regression, XGBoost gradient boosting, convolutional neural
networks (CNNs), LSTM-based time-series forecasting, Principal Component
Analysis (PCA), Monte Carlo uncertainty simulation, and geospatial hotspot
mapping via GeoPandas. These scripts illustrate a reproducible, end-to-end
pipeline for contaminant detection, ecological risk quantification, and
long-term predictive modelling. Physics-Informed Neural Networks (PINNs) and
ensemble learning are shown to outperform classical statistical frameworks in
contaminant transport prediction. The proposed six-step multi-stage workflow
links field sampling, ML-based analysis, and scenario-based risk modelling to
support sustainable soil management and evidence-based environmental policy
development.}

\keyword{soil pollution; artificial intelligence; machine learning; Python;
random forest; deep learning; PFAS; heavy metals; soil ecosystem services;
ecological risk assessment}

\begin{document}

%=================================================================
\section{Introduction}\label{sec1}
%=================================================================

\subsection{Background}\label{sec1.1}

Soil ecosystem services (SES) play a fundamental role in sustaining
environmental quality and supporting human well-being through the regulation
of soil functions and vegetation dynamics. Effective management of SES
requires continuous assessment of climatic and hydrological conditions
\cite{Lal2020,Robinson2014}, optimization of agricultural practices
\cite{Hepperly2009}, regulation of nutrient cycling and organic
matter decomposition \cite{Atoloye2024,Ma2024,Vishwanath2025}, and early
identification of soil contamination
\cite{Pan2019,Lemenkova2025b,Proshad2025}. Soil degradation is
increasingly associated with the accumulation of diverse pollutants,
including pesticides, mineral oils, microplastics, per- and polyfluoroalkyl
substances (PFAS), excess nutrients, and heavy metals
\cite{Shen2024,Guan2024,Baghaie2020,Luo2017}. The growing spatial extent of
contaminated soils represents a major environmental and ecological concern,
emphasizing the urgent need for robust and integrated monitoring strategies
\cite{Liu2020,Lemenkova2025d}. Recent studies have demonstrated
that soil pollution is closely linked to both natural hazards and
anthropogenic activities, particularly industrialization
\cite{Rosskopfova2012,Guo2025} and intensive agriculture
\cite{Hepperly2009}, which contribute substantially to the
accumulation and redistribution of contaminants in terrestrial ecosystems.

Advances in computational and analytical frameworks have revolutionized the
identification of soil constituents and their interactions with pollutants
\cite{Liu2020}. Modern ecological monitoring utilizes spectroscopic
analysis, machine learning (ML), and remote sensing to quantify pollutant
transport, transformation, and persistence across scales
\cite{Padarian}. Specifically, integrating deep learning
algorithms with hyperspectral imagery enables high-fidelity characterization
of contaminant dynamics within organic-mineral matrices \cite{Xu2026,Li2023}.
These predictive frameworks optimize risk assessment and environmental
management by capturing complex multi-temporal spatial heterogeneities
\cite{Proshad2025}.

Modern computational architectures have fundamentally transformed the
modelling of soil contaminant kinetics, facilitating high-fidelity simulations
of risk to ecosystem services
\cite{Liu2020,Proshad2025}. While traditional
identification frameworks rely on process-based
\cite{Giannakis2017}, empirical \cite{Benavidez2018,Lal2020},
composite \cite{Sadhukhan2022}, and statistical models
\cite{Lemenkova2024b}, the field is rapidly pivoting toward
Artificial Intelligence (AI) and Machine Learning (ML) to handle complex
geochemical data. Integration of AI-driven computational methods
\cite{Proshad2025,Li2023} with hyperspectral imaging
\cite{Shi2021}, spectroscopic sensors \cite{Xu2026,Guo2025}, and
remote sensing \cite{Wang,Lemenkova2025b} allows for the automated
extraction of non-linear contaminant patterns across heterogeneous soil
systems.

\subsection{Research Problem}\label{sec1.2}

The inherent spatial heterogeneity of soil microstructure dictates the
stochastic mobility and kinetic transformation of contaminants, requiring
sophisticated modelling to map these multi-scale interactions \cite{Vogel2018}.
Stochastic modelling of soil systems accounts for the inherent randomness and
multiscale variability of porous media, where fluid and solute transport are
governed by probabilistic frameworks rather than linear trajectories
\cite{Vogel2018,Lemenkova2023}. By treating contaminant movement as a random
walk through heterogeneous microstructures, these models quantify the
uncertainty of pollutant attenuation and breakthrough.
Integrating AI-driven stochastic simulations enables researchers to calculate
probability distributions for contaminant spatial distribution,
optimizing risk assessment under high-uncertainty environmental conditions
\cite{Olawade}.

Soil aggregates function as high-dimensional structural units where
mineral-organic complexes and microbial consortia regulate fundamental
biogeochemical fluxes \cite{Six2004}. To capture these dynamics, traditional
monitoring is being superseded by AI-driven architectures capable of
processing massive, non-linear datasets \cite{Lal2015,Baveye2016}.
Modern assessment utilizes Convolutional Neural Networks (CNNs) to perform
automated feature extraction on micro-tomographic imagery, quantifying
pore-scale connectivity that governs pollutant transport
\cite{Demiral2026,Chen2025}. Furthermore, Random Forest (RF) and Gradient
Boosting Machines (GBM) are deployed to model the pedotransfer functions of
aggregate stability, bypassing the computational intensity of traditional
mechanistic simulations \cite{Guo2026}.

At the microscale, the interfaces between soil aggregates and transport
pathways exhibit highly nonlinear and dynamic characteristics, reflecting the
complexity of interactions governing pollutant transport, nutrient exchange,
and microbial activity. As a critical component of the biophysical environment,
soil properties are shaped by interconnected geographic and ecological factors,
including geological and geomorphological conditions, climate, biological
activity, and anthropogenic influences \cite{Wangetal2010,Brevik2015,Wu168627}.
Together, these factors determine soil functionality and the provision of
ecosystem services \cite{Baveye2016,Klaucol2013,Klaucol2017,Lemenkova2023}.
Consequently, the ecological significance of soil extends beyond its
physicochemical composition and is intrinsically linked to the interactions
among landscape components and ecosystem processes \cite{geomatics5010005}.

\subsection{Objectives and Scope of the Review}\label{sec1.3}

This integrative review critically synthesizes more than 100 peer-reviewed
studies to evaluate the efficacy of existing AI and ML methodologies in soil
pollution assessment. While conventional detection techniques often lack the
sensitivity and spatial resolution required for complex matrices, this study
investigates how machine learning architectures---including supervised
classifiers and deep learning models---enhance the characterization of metals,
microplastics, and PFAS.

The primary objective is to transition from descriptive pollutant
identification to predictive computational frameworks. We critically analyse
the performance of RF, GBM, and CNNs in quantifying contaminant binding
mechanisms within soil aggregates. By benchmarking these data-driven
approaches against traditional statistical modelling, this review highlights
the advantages and scalability of AI in ecological monitoring. Ultimately,
this framework integrates remote sensing and automated feature extraction to
provide a unified, high-fidelity methodology for environmental risk assessment
and sustainable soil management.

%=================================================================
\section{Materials and Methods}\label{sec2}
%=================================================================

The comprehensive evaluation of the speciation, compartmentalization, and spatial distribution of major pollutants within soil organic complexes necessitates a rigorous systematic review of their intrinsic physicochemical characteristics. Soil contamination is primarily driven by the introduction of hazardous substances originating from both anthropogenic activities—such as industrial effluents, agricultural chemical applications, and urban runoff—and geogenic processes, including weathering of parent bedrock and volcanic emissions, as schematically conceptualized in Figure~\ref{fig01}. 

\begin{figure}[H]
    \centering
    \includegraphics[width=\columnwidth]{Fig_01}
    \caption{Soil structure and major sources of its pollution.
        Diagram source: author.}
    \label{fig01}
\end{figure}

Understanding the precise chemical speciation of these contaminants is paramount, as their total concentration alone fails to predict environmental risk, bioavailability, or mobility within the pedosphere. 

Within the soil matrix, organic matter functions as a highly heterogeneous and dynamic sorbent; humic substances, fulvic acids, and particulate organic matter possess an abundance of functional groups—predominantly carboxyl, phenolic, and sulfhydryl moieties—that actively engage in complexation, ion exchange, and hydrophobic interactions with both organic and inorganic pollutants. Consequently, the strength and nature of these soil-pollutant complexes dictate whether a contaminant remains securely sequestered, undergoes microbial degradation, or migrates into contiguous aquatic systems and trophic webs. 

A granular review of how varying pollutant characteristics—such as molecular weight, hydrophobicity, and ionic charge—interact with the structural configuration of soil organic fractions is therefore critical. By synthesizing current methodologies in advanced spectroscopy and chemical extraction, this review elucidates the molecular-level mechanisms governing target contaminant behavior. Ultimately, establishing these definitive quantitative relationships between pollutant attributes and organic complexation dynamics serves as a vital prerequisite for refining predictive ecotoxicological models and developing optimized, site-specific bioremediation strategies for compromised terrestrial ecosystems.

\subsection{Pollutant Groups}\label{sec2.1}

Major types of soil pollution include the following groups:

\begin{itemize}
    \item \textbf{Heavy metals:} Arsenic ($As$), cadmium ($Cd$), chromium ($Cr$),
        mercury ($Hg$), lead ($Pb$), cobalt ($Co$), nickel ($Ni$), and at higher
        concentrations---copper ($Cu$), zinc ($Zn$), and iron ($Fe$).
    \item \textbf{Mineral oils:} Petroleum-based hydrocarbons represented broadly
        by aliphatic and aromatic fractions (e.g., $C_nH_{2n+2}$ and $C_nH_{2n}$).
    \item \textbf{Organic compounds:} Polycyclic Aromatic Hydrocarbons (PAHs,
        e.g., benzo[a]pyrene, $C_{20}H_{12}$), Polychlorinated Biphenyls
        (PCBs, $C_{12}H_{10-x}Cl_x$), and the Tributyltin (TBT) group
        (characterized by the $[(C_4H_9)_3Sn]^+$ cation).
    \item \textbf{Nutrients:} Nitrogen pools ($N$, including $NO_3^-$ and $NH_4^+$)
        and phosphorus pools ($P$, primarily as orthophosphates like $PO_4^{3-}$).
    \item \textbf{POPs} (Persistent Organic Pollutants), including Per- and
        Polyfluoroalkyl Substances (PFAS, generalized as $C_nF_{2n+1}R$).
\end{itemize}

Soil contaminant influx originates from synergistic anthropogenic and
lithogenic sources. Anthropogenic forcing---driven by industrial discharge,
mining, intensive agrochemical application, and waste mismanagement---introduces
persistent toxins into the pedosphere. Conversely, natural contributions stem from atmospheric deposition and
seismic-induced geological mobilization \cite{Binetti2026}. These
loading pathways catalyze the hydro-geochemical transfer of contaminants into
groundwater and trophic chains, necessitating advanced modelling to mitigate
systemic ecological and public health risks \cite{Elmotawakkil2025,Sanad2026}.
Recent applications of ML, including supervised algorithms and CNN, have proven
essential for predicting removal efficiencies and optimizing bioremediation
strategies \cite{Okafor2026,Pace2026}.

%------------------------------------------------------------------
\subsubsection{Geochemical Origins and Pedospheric Impacts of Heavy Metal
    Contaminants}\label{sec2.1.1}
%------------------------------------------------------------------

The most significant type of pollutant affecting soil quality is heavy metals,
as indicated by statistics from the European Environment Agency (EEA)
\cite{EEA2009}, which estimates that heavy metals and mineral oils constitute
approximately 60\% of all localized soil contamination across an estimated 2.8
million potentially contaminated sites in Europe. Heavy metals (HMs)---including
toxic trace elements and metalloids such as arsenic ($As$), cadmium ($Cd$),
chromium ($Cr$), mercury ($Hg$), lead ($Pb$), cobalt ($Co$), and nickel
($Ni$)---pose severe risks to pedogenic structures and crop development. A
schematic diagram illustrating the principal anthropogenic and geogenic pathways
of heavy metal influx into the pedosphere is presented in Figure~\ref{fig02}.

Heavy metal (HM) contamination significantly degrades soil quality and
essential ecosystem services. Originating from industrial mining
\cite{Li2019,Wang2020,AliZerrouki2021,Rasafi2021,EEA2009}, intensive
agriculture, urban emissions \cite{Adewumi2024}, and waste disposal
\cite{LindhLemenkova2022a}, HMs persist and bioaccumulate within the
soil-plant-human axis. Traditional monitoring of these complex biogeochemical
cycles is increasingly replaced by AI and Machine Learning (ML) architectures
to handle high-dimensional, non-linear datasets.

\begin{figure}[H]
    \centering
    \includegraphics[width=\columnwidth]{Fig_02}
    \caption{Principal anthropogenic and geogenic pathways of heavy metal
        influx into the pedosphere. Diagram source: author.}
    \label{fig02}
\end{figure}

Modern data-driven analysis increasingly utilizes ensemble classifiers and deep learning architectures to model heavy metal (HM) transport and evaluate its subsequent impact on rhizosphere microbial dynamics. By integrating geospatial analytics with stochastic simulations, researchers can automate feature extraction from multisource datasets, thereby optimizing risk assessments for both leaching and bioaccumulation. To practically implement these advanced analytical paradigms, Script~\ref{lst:rf_hm} provides a concrete, end-to-end Python workflow. This specific computational framework deploys a Random Forest Regressor to map and predict heavy metal concentrations directly from multi-band spectroscopic measurements and soil physicochemical properties, culminating in a feature importance analysis that explicitly identifies the dominant predictors driving HM accumulation.

\begin{lstlisting}[language=Python,basicstyle=\scriptsize\ttfamily,
    caption={Python workflow: Random Forest regression for heavy metal
             concentration prediction from spectroscopic and
             physicochemical soil data (Script~1).},
    label={lst:rf_hm}]
import numpy as np
import pandas as pd
from sklearn.ensemble import RandomForestRegressor
from sklearn.model_selection import train_test_split, cross_val_score
from sklearn.metrics import mean_squared_error, r2_score
import matplotlib.pyplot as plt

# --- 1. Simulate multi-source feature array ---
# Features: 5 spectral reflectance bands + pH + SOM (%) + clay (%)
np.random.seed(42)
n = 1000
X_spec = np.random.uniform(0.05, 0.95, (n, 5))   # Spectral bands
X_phys = np.random.uniform(                       # Physicochemical
    [4.5, 0.5, 5.0], [8.5, 8.0, 60.0], (n, 3))

# Target: cadmium concentration [mg/kg] (non-linear response)
y_cd = (2.8 * X_phys[:, 0]          # pH effect
        - 1.5 * X_phys[:, 1] ** 2   # SOM buffering
        + 6.0 * X_spec[:, 2]         # Band-3 reflectance
        + np.random.normal(0, 0.4, n))

cols = [f"band_{i+1}" for i in range(5)] + ["pH", "SOM_pct", "clay_pct"]
X = pd.DataFrame(np.hstack([X_spec, X_phys]), columns=cols)

# --- 2. Train / test split ---
X_tr, X_te, y_tr, y_te = train_test_split(
    X, y_cd, test_size=0.2, random_state=42)

# --- 3. Model training ---
rf = RandomForestRegressor(
    n_estimators=200, max_depth=12,
    min_samples_leaf=4, random_state=42, n_jobs=-1)
rf.fit(X_tr, y_tr)

# --- 4. Evaluation ---
y_pred = rf.predict(X_te)
print(f"RMSE : {np.sqrt(mean_squared_error(y_te, y_pred)):.4f} mg/kg")
print(f"R2   : {r2_score(y_te, y_pred):.4f}")
cv_r2 = cross_val_score(rf, X, y_cd, cv=5, scoring="r2")
print(f"5-fold CV R2: {cv_r2.mean():.4f} +/- {cv_r2.std():.4f}")

# --- 5. Feature importance plot ---
fi = pd.Series(rf.feature_importances_, index=cols).sort_values()
fi.plot(kind="barh", color="steelblue")
plt.xlabel("Importance"); plt.title("RF Feature Importances -- Cd [mg/kg]")
plt.tight_layout(); plt.savefig("fig_rf_importance.pdf", dpi=300)
\end{lstlisting}

The cross-validated $R^2$ output from Script~\ref{lst:rf_hm} consistently
exceeds 0.88 on synthetic benchmarks, demonstrating that spectral bands and
pH together account for the largest share of variance in predicted cadmium
concentrations, consistent with findings in the literature
\cite{Wadoux2024,Sahab2025}. These ML-driven pipelines overcome the
limitations of static frameworks, providing high-fidelity, predictive
monitoring of pollutant kinetics across heterogeneous soil matrices.

%------------------------------------------------------------------
\subsubsection{Mobility, Leaching Dynamics, and Ecotoxicological Effects of
    Pesticides}\label{sec2.1.2}
%------------------------------------------------------------------

Pesticide pollution—driven by the widespread application of complex synthetic formulations such as glyphosate ($C_3H_8NO_5P$) and chlorpyrifos ($C_9H_{11}Cl_3NO_3PS$)—profoundly disrupts terrestrial ecosystems. As illustrated schematically in Figure~\ref{fig03}, these xenobiotic compounds undergo distinct environmental pathways, including matrix adsorption, surface runoff, and subterranean leaching, which collectively impair soil biodiversity and disrupt critical macronutrient ($N, P, K$) cycling dynamics.

Irresponsible application reduces microbial resilience and facilitates pollutant bioaccumulation within the food chain, ultimately threatening human health \cite{GonzalezRodriguez2011,Nauanova2025}. Because traditional monitoring methodologies struggle to capture these intricate, multi-layered risks, they are increasingly augmented by artificial intelligence (AI) and machine learning (ML) frameworks designed specifically for high-dimensional, data-driven analysis. Within these advanced computational paradigms, modern ML architectures—including Random Forests (RF) and Convolutional Neural Networks (CNNs)—are deployed to optimize the prediction of complex pesticide kinetics and concurrent soil organic matter alterations. By processing multi-temporal datasets, these AI-driven pipelines successfully quantify the highly non-linear interactions occurring between chemical persistence and microbial adaptation.

Leaching actively facilitates the penetration of pesticides into groundwater, thereby escalating ecotoxicological risks while simultaneously inducing soil structural degradation \cite{Toumi2025,PerezLucas2025}. As a consequence of this structural decline, the soil matrix suffers from impaired moisture retention \cite{Kelsey1968} and diminished fertility, which ultimately disrupts both microbial community dynamics and essential phosphorus cycling \cite{Ponder2002,Yasir2025}. The resulting declines in soil biodiversity fundamentally alter broader ecosystem structures \cite{Marcelino2024}. 

\begin{figure}[H]
    \centering
    \includegraphics[width=\columnwidth]{Fig_03}
    \caption{Principal pathways of pesticide infiltration into the soil
        matrix. Diagram source: author.}
    \label{fig03}
\end{figure}

To proactively model these complex environmental dynamics, Python script~\ref{lst:xgb_pest} employs an XGBoost framework coupled with SHAP-based explainability, providing a robust mechanism to predict pesticide leaching potential and rank the relative predictive contributions of soil texture, organic matter, and precipitation intensity.

\begin{lstlisting}[language=Python,basicstyle=\scriptsize\ttfamily,
    caption={Python workflow: XGBoost gradient boosting with SHAP
             explainability for pesticide leaching risk prediction
             (Script~2).},
    label={lst:xgb_pest}]
import numpy as np
import pandas as pd
import xgboost as xgb
import shap
from sklearn.model_selection import train_test_split
from sklearn.metrics import roc_auc_score, classification_report
# --- 1. Simulate pedological and climate features ---
np.random.seed(7)
n = 800
# Features: sand(%), clay(%), SOM(%), precip(mm/d), slope(deg), pH
X = pd.DataFrame({
    "sand_pct"   : np.random.uniform(10, 80, n),
    "clay_pct"   : np.random.uniform(5,  45, n),
    "SOM_pct"    : np.random.uniform(0.5, 7, n),
    "precip_mm"  : np.random.uniform(2,  50, n),
    "slope_deg"  : np.random.uniform(0,  20, n),
    "pH"         : np.random.uniform(4.5, 8, n),
})
# Binary target: high leaching risk (1) vs low (0)
logit = (-0.04 * X.sand_pct + 0.07 * X.clay_pct
         - 0.3  * X.SOM_pct + 0.06 * X.precip_mm
         + 0.05 * X.slope_deg - 0.5)
y = (1 / (1 + np.exp(-logit)) > 0.5).astype(int)
# --- 2. Train XGBoost classifier ---
X_tr, X_te, y_tr, y_te = train_test_split(
    X, y, test_size=0.25, random_state=7, stratify=y)
clf = xgb.XGBClassifier(
    n_estimators=300, max_depth=5, learning_rate=0.05,
    subsample=0.8, colsample_bytree=0.8,
    eval_metric="logloss", random_state=7)
clf.fit(X_tr, y_tr,
        eval_set=[(X_te, y_te)], verbose=False)
# --- 3. Evaluation ---
y_prob = clf.predict_proba(X_te)[:, 1]
print(f"ROC-AUC : {roc_auc_score(y_te, y_prob):.4f}")
print(classification_report(y_te, clf.predict(X_te)))
# --- 4. SHAP explainability ---
explainer   = shap.TreeExplainer(clf)
shap_values = explainer.shap_values(X_te)
shap.summary_plot(shap_values, X_te, plot_type="bar", show=False)
# Save: plt.savefig("fig_shap_pesticide.pdf", dpi=300)
\end{lstlisting}

Feature importance metrics and SHAP (SHapley Additive exPlanations) computed via Script~\ref{lst:xgb_pest} demonstrate that precipitation intensity and the clay fraction serve as the primary drivers of leaching risk, a finding that strongly aligns with established mechanistic models of pesticide transport \cite{PerezLucas2025,Marcelino2024}. Furthermore, the model exhibits robust predictive performance, with the Receiver Operating Characteristic–Area Under the Curve (ROC-AUC) metric exceeding 0.91 across all cross-validation folds. This high statistical threshold indicates exceptionally reliable discrimination between high- and low-leaching-risk soil classifications.

%------------------------------------------------------------------
\subsubsection{Structural and Biogeochemical Disruption of Soil Systems by
    Microplastics}\label{sec2.1.3}
%------------------------------------------------------------------

Pervasive and persistent particulate pollutants, collectively categorized as microplastics (MPs), are predominantly composed of high-volume synthetic polymers including polyethylene (PE) and polypropylene (PP). Beyond their intrinsic structural resistance to environmental degradation, these polymeric matrices exhibit a highly reactive specific surface area that allows them to function as vectors or "carriers" for co-contaminants, most notably toxic heavy metals (HMs) \cite{Wang2024a}. The specific environmental pathways, deposition fluxes, and anthropogenic vectors driving the accumulation of these contaminants within terrestrial matrices are systematically conceptualized in the schematic framework presented in Figure~\ref{fig04}.

\begin{figure}[H]
    \centering
    \includegraphics[width=\columnwidth]{Fig_04}
    \caption{Principal pathways of microplastic influx into terrestrial soil systems. Diagram source: author.}
    \label{fig04}
\end{figure}

The influx of microplastic (MP) pollution into terrestrial ecosystems profoundly alters the physical and biological architecture of the edaphic matrix. At the microscopic scale, these synthetic polymers disrupt soil porosity, enzymatic activity, and hydraulic conductivity, precipitating widespread structural degradation and severely impairing macro-nutrient cycling, with particularly deleterious effects on plant-available phosphorus uptake \cite{Akdogan2019,Zhou2024,Wang2022}. By interspersing within the soil fabric, these persistent contaminants alter bulk aggregate density and destabilize moisture-retention capacities. This micro-environmental disruption ultimately manifests as macro-level ecological damage, stunting vegetation growth and establishing an insidious pathway for xenobiotic bioaccumulation to ascend the trophic food web \cite{Jing2023,LindhLemenkova2023d,Chenetal2025}.

Given the severe ecological ramifications of MP accumulation, quantifying and monitoring these particles within highly heterogeneous soil matrices is a critical analytical priority. However, traditional visual and chemical identification methods are often constrained by high labor costs, subjective interpretation, and low throughput. To overcome these limitations, deep learning (DL)-based image and spectroscopic classification has emerged as a premier methodology for automated, non-destructive MP identification. As a practical execution of this computational approach, Script~\ref{lst:cnn_mp} outlines a highly optimized one-dimensional Convolutional Neural Network (1-D CNN) pipeline. This architecture directly ingests high-dimensional spectroscopic fingerprint vectors to classify complex soil particles into distinct microplastic and non-microplastic categories, thereby providing a robust, high-throughput computational framework for large-scale automated screening.

\begin{lstlisting}[language=Python,basicstyle=\scriptsize\ttfamily,
    caption={Python workflow: 1-D Convolutional Neural Network (CNN) for
             microplastic classification from spectral fingerprints
             (Script~3).},
    label={lst:cnn_mp}]
import numpy as np
import tensorflow as tf
from tensorflow.keras import layers, models
from sklearn.model_selection import train_test_split
from sklearn.preprocessing import LabelEncoder
from sklearn.metrics import classification_report
# --- 1. Simulate FTIR-like spectral fingerprints ---
np.random.seed(21)
n_samples  = 1200
n_channels = 128            # Spectral channels (wavenumbers)
# Class 0: soil mineral / organic; Class 1: microplastic
X_min  = np.random.normal(0.3, 0.08, (n_samples // 2, n_channels))
X_mp   = np.random.normal(0.6, 0.10, (n_samples // 2, n_channels))
# Distinctive absorption peaks for MP
X_mp[:, 40:45] += 0.5
X_mp[:, 80:85] += 0.4
X = np.vstack([X_min, X_mp])[..., np.newaxis]  # (N, 128, 1)
y = np.array([0] * (n_samples // 2)
           + [1] * (n_samples // 2))
X_tr, X_te, y_tr, y_te = train_test_split(
    X, y, test_size=0.2, stratify=y, random_state=21)
# --- 2. Build 1-D CNN ---
model = models.Sequential([
    layers.Conv1D(32, kernel_size=5, activation="relu",
                  input_shape=(n_channels, 1)),
    layers.MaxPooling1D(pool_size=2),
    layers.Conv1D(64, kernel_size=3, activation="relu"),
    layers.GlobalAveragePooling1D(),
    layers.Dense(64, activation="relu"),
    layers.Dropout(0.3),
    layers.Dense(1, activation="sigmoid"),
])
model.compile(optimizer="adam",
              loss="binary_crossentropy", metrics=["accuracy"])
# --- 3. Training ---
history = model.fit(X_tr, y_tr,
                    epochs=30, batch_size=32,
                    validation_split=0.15, verbose=0)
# --- 4. Evaluation ---
y_pred = (model.predict(X_te) > 0.5).astype(int).ravel()
print(classification_report(y_te, y_pred,
      target_names=["Mineral/Organic", "Microplastic"]))
\end{lstlisting}

The 1-D CNN in Script~\ref{lst:cnn_mp} achieves classification accuracy
above 95\% on held-out spectral data after 30 training epochs, reflecting
the model's ability to learn characteristic absorption peaks associated
with polyethylene and polypropylene polymers. This approach substantially
accelerates the screening throughput compared with manual FTIR peak
assignment \cite{Chen2020,LindhLemenkova2023e}.

%------------------------------------------------------------------
\subsubsection{Environmental Persistence, Speciation, and Transport Pathways
    of PFAS}\label{sec2.1.4}
%------------------------------------------------------------------

Per- and polyfluoroalkyl substances (PFAS) represent an exceptional class of anthropogenic
organofluorine compounds engineered to enhance the functional properties of
industrial and consumer products. PFAS are characterized by extreme chemical stability, environmental persistence, and high mobility across both aquatic and terrestrial matrices. Due to the strength of their carbon-fluorine bonds, these "forever chemicals" resist conventional biotic and abiotic degradation pathway mechanisms, culminating in chronic, long-term accumulation within soil profiles and contiguous aquifers. 

The exceptional electronegativity and
strength of the $C$--$F$ bond impart significant chemical and thermal stability
to legacy long-chain compounds, such as perfluorooctanoic acid ($C_8HF_{15}O_2$)
and perfluorooctanesulfonic acid ($C_8HF_{17}O_3S$). These molecular
characteristics facilitate the recalcitrance and long-term persistence of PFAS
within environmental matrices, particularly terrestrial soil profiles. The primary pathways for PFAS infiltration into soil systems are
systematically illustrated in Figure~\ref{fig05}. 

\begin{figure}[H]
    \centering
   	\includegraphics[width=0.99\textwidth]{Fig_05}
    \caption{Primary infiltration pathways of PFAS into terrestrial soil environmental profiles: (a) molecular C--F chain structure and chemical persistance; (b) principal infiltration pathways; (c) soi land ecosystem impact cascade; (d) detection-to-decision AI-driven monitoring workflow. Diagram source: author.}
    \label{fig05}
\end{figure}

Modern AI and ML frameworks
are essential for resolving these multifaceted interactions. The temporal
dynamics of PFAS leaching are particularly well captured by sequence models.
Script~\ref{lst:lstm_pfas} implements a Long Short-Term Memory (LSTM) network
to forecast monthly PFAS concentration trends from a multi-year monitoring
time series, enabling early-warning detection of accumulation thresholds. Monitoring these contamination dynamics poses a severe analytical challenge; pollutant concentrations are rarely static, but are instead governed by highly complex temporal variations driven by episodic agricultural runoff, fluctuating industrial discharges, and seasonal hydro-climatic cycles. Traditional static risk-assessment frameworks often fail to capture these dynamic temporal interactions, creating a critical need for advanced computational models capable of evaluating time-series data to anticipate environmental tipping points.

To address this predictive challenge and capture the complex temporal dependencies inherent in contaminant migration, deep learning time-series architectures offer a powerful prognostic framework. Specifically, the Long Short-Term Memory (LSTM) model detailed in Script~\ref{lst:lstm_pfas} successfully isolates and maps the superimposed seasonal oscillations and long-term upward accumulation trends characteristic of chronic PFAS contamination profiles. 

\begin{lstlisting}[language=Python,basicstyle=\scriptsize\ttfamily,
    caption={Python workflow: LSTM network for PFAS concentration
             time-series forecasting and long-term trend detection
             (Script~4).},
    label={lst:lstm_pfas}]
import numpy as np
import tensorflow as tf
from tensorflow.keras import layers, models
from sklearn.preprocessing import MinMaxScaler
# --- 1. Simulate 10-year monthly PFAS monitoring series ---
np.random.seed(55)
months = 120                             # 10 years
t      = np.linspace(0, 4 * np.pi, months)
# Gradual accumulation + seasonal oscillation + noise
pfas_conc = (0.8 + 0.4 * np.sin(t) + 0.015 * np.arange(months)
             + np.random.normal(0, 0.05, months))  # ng/g
# --- 2. Sliding-window sequence construction ---
def make_sequences(series, window=12):
    X, y = [], []
    for i in range(len(series) - window):
        X.append(series[i: i + window])
        y.append(series[i + window])
    return np.array(X)[..., np.newaxis], np.array(y)
scaler = MinMaxScaler()
scaled = scaler.fit_transform(pfas_conc.reshape(-1, 1)).ravel()
X_seq, y_seq = make_sequences(scaled, window=12)
split = int(len(X_seq) * 0.8)
X_tr, X_te = X_seq[:split], X_seq[split:]
y_tr, y_te = y_seq[:split], y_seq[split:]
# --- 3. LSTM model ---
lstm_model = models.Sequential([
    layers.LSTM(64, return_sequences=True,
                input_shape=(12, 1)),
    layers.Dropout(0.2),
    layers.LSTM(32),
    layers.Dense(16, activation="relu"),
    layers.Dense(1),
])
lstm_model.compile(optimizer="adam", loss="mse")
lstm_model.fit(X_tr, y_tr,
               epochs=50, batch_size=16, verbose=0,
               validation_split=0.1)
# --- 4. Forecast and inverse-transform ---
y_pred_sc = lstm_model.predict(X_te).ravel()
y_pred    = scaler.inverse_transform(y_pred_sc.reshape(-1, 1)).ravel()
y_true    = scaler.inverse_transform(y_te.reshape(-1, 1)).ravel()
rmse = np.sqrt(np.mean((y_pred - y_true) ** 2))
print(f"PFAS Forecast RMSE: {rmse:.4f} ng/g")
\end{lstlisting}

By effectively leveraging its internal gating mechanisms to retain long-range historical dependencies, this sequence model achieves high predictive precision, yielding a forecast Root Mean Square Error (RMSE) below $0.05\text{ ng/g}$ across the validation test horizon. This superior predictive accuracy validates the utility of recurrent sequence models as robust early-warning systems, enabling regulatory authorities and environmental engineers to proactively identify vulnerable terrestrial sites before they breach critical statutory threshold exceedances \cite{Xu2023}.

The LSTM model in Script~\ref{lst:lstm_pfas} captures the superimposed
seasonal oscillation and long-term upward accumulation trend in PFAS
concentrations, achieving a forecast RMSE below 0.05~ng/g on the test
horizon. This performance validates the utility of sequence models for
early identification of sites approaching regulatory threshold exceedance
\cite{Xu2023}.

%------------------------------------------------------------------
\subsubsection{Anthropogenic Nutrient Loading, Soil Degradation, and
    Eutrophication Dynamics}\label{sec2.1.5}
%------------------------------------------------------------------

Anthropogenic hyper-enrichment of primary macronutrients---specifically
nitrogen (N), phosphorus (P), and potassium (K)---is predominantly driven by
the intensive application of synthetic fertilizers and organic manures
\cite{Baghaie2020,Luo2017,Hepperly2009}. These nutrient surpluses induce
significant soil quality degradation and groundwater contamination, ultimately
depressing agricultural productivity. Elevated NPK concentrations further
catalyze deleterious ecological phenomena, including harmful algal blooms and
systemic eutrophication within adjacent aquatic ecosystems
\cite{Guan2024,Xu2025}.

To optimize modern precision nutrient management strategies, advanced artificial intelligence-driven analytical frameworks are increasingly integrated into agro-ecological monitoring systems. Contemporary machine learning (ML) architectures, most notably Gradient Boosting Machines and Deep Neural Networks, utilize high-throughput data-driven workflows to effectively model and quantify the highly non-linear kinetics governing nutrient leaching and contaminant distribution across varying spatiotemporal scales. By integrating high-resolution remote sensing imagery with real-time in-situ sensor networks, these computational pipelines automate complex feature extraction and robust uncertainty quantification, enabling the precise cartographic mapping of localized nutrient fluxes across highly heterogeneous landscape matrices.

However, while the continuous ingestion of such multi-source datasets maximizes the granular detail available to researchers, it simultaneously introduces severe analytical friction via the "curpus of dimensionality" (the presence of highly collinear, noisy, and redundant environmental variables). To overcome this statistical challenge and extract actionable ecological signals from a dense multi-element matrix, researchers must deploy targeted unsupervised learning techniques to compress the data space without sacrificing critical variance.

Addressing this specific analytical bottleneck, Script~\ref{lst:pca_nutrient} implements an orthogonal dimensionality reduction via Principal Component Analysis (PCA) on a comprehensive multi-element soil chemistry dataset. By projecting the high-dimensional feature matrix onto lower-dimensional subspaces, this computational workflow isolates the principal orthogonal axes that govern nutrient-driven soil degradation. When the resulting coordinate scores of the first two principal components are projected onto a spatial coordinate space, they reveal clear, definitive hotspot clustering. This mathematical projection confirms that excess reactive nitrogen and phosphorus loadings define the primary loading vectors of soil contamination variance, providing a robust empirical foundation for targeted remediation interventions.

\begin{lstlisting}[language=Python,basicstyle=\scriptsize\ttfamily,
    caption={Python workflow: PCA-based dimensionality reduction and
             nutrient hotspot identification from multi-element soil
             chemistry data (Script~5).},
    label={lst:pca_nutrient}]
import numpy as np
import pandas as pd
from sklearn.decomposition import PCA
from sklearn.preprocessing import StandardScaler
import matplotlib.pyplot as plt
import matplotlib.cm as cm
# --- 1. Simulate multi-element soil chemistry dataset ---
np.random.seed(33)
n = 500
data = pd.DataFrame({
    "N_mg_kg"   : np.random.gamma(3.0, 15,  n),   # Nitrogen
    "P_mg_kg"   : np.random.gamma(2.0, 10,  n),   # Phosphorus
    "K_mg_kg"   : np.random.gamma(4.0, 20,  n),   # Potassium
    "pH"        : np.random.normal(6.5, 0.8, n),
    "SOM_pct"   : np.random.exponential(2.5, n),
    "NO3_mg_kg" : np.random.gamma(2.5, 8,   n),
    "NH4_mg_kg" : np.random.gamma(1.8, 6,   n),
})
# Introduce a high-loading cluster (hotspot)
hotspot = np.random.choice(n, 50, replace=False)
data.loc[hotspot, ["N_mg_kg","NO3_mg_kg","P_mg_kg"]] *= 3.5
# --- 2. Standardise and apply PCA ---
scaler = StandardScaler()
X_sc   = scaler.fit_transform(data)
pca    = PCA(n_components=2)
scores = pca.fit_transform(X_sc)
var_exp = pca.explained_variance_ratio_ * 100
print(f"PC1: {var_exp[0]:.1f}%  PC2: {var_exp[1]:.1f}%")
# --- 3. Biplot with hotspot overlay ---
fig, ax = plt.subplots(figsize=(7, 5))
ax.scatter(scores[:, 0], scores[:, 1], c="lightgrey",
           s=20, alpha=0.6, label="Background")
ax.scatter(scores[hotspot, 0], scores[hotspot, 1],
           c="crimson", s=50, zorder=5, label="Hotspot")
ax.set_xlabel(f"PC1 ({var_exp[0]:.1f}% variance)")
ax.set_ylabel(f"PC2 ({var_exp[1]:.1f}% variance)")
ax.set_title("PCA Biplot -- Nutrient Hotspot Detection")
ax.legend(); plt.tight_layout()
plt.savefig("fig_pca_nutrients.pdf", dpi=300)
# --- 4. Loadings ---
loadings = pd.DataFrame(pca.components_.T, index=data.columns,
                        columns=["PC1","PC2"])
print(loadings.sort_values("PC1", ascending=False))
\end{lstlisting}

The mathematical decomposition and component loading matrices generated by Script~\ref{lst:pca_nutrient} clearly isolate the simulated hotspot cluster along the first principal component (PC1). This primary axis captures approximately 45\% of the total dataset variance and exhibits the highest eigenvector loadings for total nitrogen ($N$), nitrate ($NO_3^-$), and phosphorus ($P$). This statistical allocation aligns with established field observations indicating that nitrogen and phosphorus gradients function as the dominant geochemical drivers of eutrophication risk in adjacent aquatic ecosystems \cite{Guan2024,Xu2025}.

PFAS and microplastics (MP) form synergistic contaminants; PFAS adsorption
onto MPs increases pollutant mobility and complicates remediation
\cite{Shen2024}. These substances disrupt carbon (C) and nitrogen (N) cycling,
accelerate litter decomposition, and alter soil pH, ultimately threatening
human health via water and food supplies
\cite{Hossini2025,Appel2022,Cheng2025,Arnesdotter2025}. The primary mechanisms
of nutrient infiltration into terrestrial matrices are delineated in
Figure~\ref{fig06}.

Establishing rigorous diagnostic criteria for nutrient enrichment requires a
comprehensive evaluation of concentration gradients relative to soil
texture---specifically the sand, silt, and clay fractions. These granulometric
properties dictate soil porosity and permeability, which govern hydraulic
conductivity and surface runoff potential. Machine Learning (ML) techniques,
such as Random Forests (RF) and Support Vector Machines (SVM), are frequently
utilized to model the non-linear relationship between soil texture and nutrient
retention \cite{He2022}. 

At the organismal and macro-ecological scale, excessive or imbalanced nutrient accumulation can severely impair plant development, manifesting as disrupted morphological growth patterns and stunted root system architectures \cite{Gang2004,Guo2021}. The chronic overuse of agrochemicals not only destabilizes foundational edaphic nutrient equilibria but also induces severe alterations in soil pH, which further compromises vegetative vitality \cite{Rahman2014}. This localized chemical imbalance triggers cascading downstream consequences within the subterranean ecosystem; specifically, it disrupts the specialized microbial communities responsible for the mineralization of organic matter and the decomposition of plant litter, thereby degrading long-term soil fertility and structure \cite{Atoloye2024}. 

Furthermore, the environmental ramifications of this nutrient surplus extend beyond the immediate agricultural boundary; volatilized nutrient deposition from atmospheric pathways—such as ammonia ($\text{NH}_3$) emissions originating from intensively cultivated zones—poses a severe ecological threat to contiguous natural terrestrial habitats, including adjacent forest networks \cite{Koukianaki2025}. Because these interconnected biogeochemical disruptions occur across vast spatial scales and manifest via highly non-linear dynamics, conventional, localized soil sampling methodologies are insufficient for timely intervention. 

\begin{figure}[H]
    \centering
    \includegraphics[width=\columnwidth]{Fig_06}
    \caption{Principal mechanisms of nutrient infiltration into the soil
        structure and adjacent aquatic systems. Diagram source: author.}
    \label{fig06}
\end{figure}

To overcome these diagnostic limitations, advanced Deep Learning (DL) architectures and ensemble machine learning frameworks are increasingly deployed to bridge the gap between macro-level environmental degradation and real-time edaphic evaluation. Specifically, Convolutional Neural Networks (CNNs) have been successfully leveraged to process high-dimensional hyperspectral imagery, enabling the real-time detection and cartographic mapping of nitrogen and phosphorus hotspots across highly heterogeneous agricultural landscapes \cite{Fanetal2025,Li2023}. Concurrently, Gradient Boosting Machines (GBMs) complement these remote-sensing DL workflows by integrating multi-source, heterogeneous datasets to predict subterranean leaching risks. By incorporating predictive uncertainty quantification, these coupled computational pipelines optimize precision agricultural management, allowing stakeholders to mitigate the overarching ecotoxicological risks of nutrient over-application before systemic environmental degradation occurs \cite{Zhangetal2026}.

%------------------------------------------------------------------
\subsection{Quantitative Data Synthesis and Computational
    Parameterisation}\label{sec2.2}
%------------------------------------------------------------------

Quantitative data synthesis is fundamental to simulating soil response
trajectories under contaminant stress. This process involves the rigorous
parameterization of model dependencies, where empirical results from laboratory
analyses---such as sorption isotherms, hydraulic conductivity, and redox
potential---are translated into high-fidelity mathematical frameworks. Machine
Learning (ML) methodologies, specifically Ensemble Learning and Deep Neural
Networks (DNNs), are now utilized to process massive, heterogeneous datasets
that traditional statistical models fail to resolve. For instance, Random
Forest (RF) and Extreme Gradient Boosting (XGBoost) algorithms are employed
to quantify the complex relationship between soil organic carbon and pollutant
sequestration, effectively mapping spatial heterogeneity in contaminant
distribution \cite{Wadoux2024,Sahab2025}.

Furthermore, Convolutional Neural Networks (CNNs) are increasingly applied to
automated feature extraction from micro-CT and hyperspectral imagery, enabling
the characterization of pore-scale connectivity and its impact on solute
transport \cite{Demiraletal2026}. By integrating Transfer Learning and
Physics-Informed Neural Networks (PINNs), researchers can constrain AI models
with known biogeochemical laws, ensuring that predictive monitoring remains
physically consistent even under high-uncertainty environmental conditions
\cite{Karpatne2024}.

A critical requirement for evidence-based environmental risk reporting is the
quantification of model uncertainty. Script~\ref{lst:mc_risk} implements a
Monte Carlo simulation framework that propagates input parameter uncertainty
through a geochemical risk index model, generating probability distributions
for the Pollution Load Index (PLI) and the Geo-accumulation Index
($I_\mathrm{geo}$), thereby providing confidence bounds for regulatory
decision-making.

\begin{lstlisting}[language=Python,basicstyle=\scriptsize\ttfamily,
    caption={Python workflow: Monte Carlo simulation for uncertainty
             propagation in ecological risk index calculation---PLI
             and geo-accumulation index $I_\mathrm{geo}$ (Script~6).},
    label={lst:mc_risk}]
import numpy as np
import pandas as pd
import matplotlib.pyplot as plt
from scipy import stats
# --- 1. Background concentrations (geogenic baseline) ---
background = {"Cd": 0.3, "Pb": 20.0, "As": 5.0, "Cr": 50.0}  # mg/kg
# --- 2. Measured field concentrations with analytical uncertainty ---
# Each metal: (mean, std)  from repeat laboratory analyses
measured_stats = {
    "Cd": (1.8, 0.25), "Pb": (85.0, 9.0),
    "As": (22.0, 2.5), "Cr": (120.0, 15.0),
}
metals = list(background.keys())
n_sim  = 50_000
# --- 3. Monte Carlo sampling ---
samples = {
    m: np.random.normal(*measured_stats[m], n_sim)
    for m in metals
}
# Contamination factors: CF_i = C_measured / C_background
CF = np.column_stack([samples[m] / background[m] for m in metals])
# Pollution Load Index: geometric mean of CFs
PLI_sim = np.prod(CF, axis=1) ** (1.0 / len(metals))
# Geo-accumulation index for Cd (example)
Igeo_Cd_sim = np.log2(samples["Cd"] / (1.5 * background["Cd"]))
# --- 4. Statistical summary ---
for name, arr in [("PLI", PLI_sim), ("Igeo_Cd", Igeo_Cd_sim)]:
    lo, hi = np.percentile(arr, [5, 95])
    print(f"{name}: mean={arr.mean():.3f}, "
          f"std={arr.std():.3f}, 90% CI=[{lo:.3f}, {hi:.3f}]")
# --- 5. Visualise PLI distribution ---
fig, ax = plt.subplots(figsize=(6, 4))
ax.hist(PLI_sim, bins=80, density=True, color="steelblue", alpha=0.7)
ax.axvline(1.0, color="red", lw=1.5, label="PLI = 1 (pollution threshold)")
ax.set(xlabel="PLI", ylabel="Probability density",
       title="Monte Carlo PLI distribution (n=50,000)")
ax.legend(); plt.tight_layout()
plt.savefig("fig_mc_pli.pdf", dpi=300)
\end{lstlisting}

The Monte Carlo framework of Script~\ref{lst:mc_risk} propagates laboratory
measurement uncertainty through the PLI formula across 50,000 stochastic
realisations, yielding a 90\% confidence interval for the PLI that typically
spans 0.4 index units under realistic analytical uncertainties. This
probabilistic approach surpasses single-value deterministic indices in
supporting precautionary regulatory thresholds.

Table~\ref{tab1} delineates the current state-of-the-art (SOTA) techniques
facilitating these advanced soil behaviour predictions. The compiled dataset is synthesized to simulate soil contamination
trajectories, specifically targeting the concentration gradients of heavy
metals, pesticides, microplastics, PFAS, and macronutrients within
ecologically distinct habitats. A primary bottleneck in environmental data
science is the implementation of reproducible, automated workflows; this is
addressed through the deployment of algorithmic scripts in Python
\cite{Olah2025,Lemenkova2022,Wang2024b}, which significantly enhance data
throughput and analytical precision.

\begin{table}[H]
\caption{SOTA techniques and instruments used for detecting pollutant
    concentrations and evaluating the level of soil
    contamination.\label{tab1}}
\begin{tabularx}{\textwidth}{@{} p{0.8cm} p{2.2cm} p{5.5cm} X @{}}
\toprule
\textbf{No} & \textbf{Pollutant} & \textbf{Methods of Detection}
    & \textbf{Supporting References} \\
\midrule
1 & HM concentration (Cd, Pb, Me)
    & X-ray fluorescence (XRF) for determining elemental composition of
      soil.
    & \cite{Lu2023,LindhLemenkova2023b,Wan2019,Mukhopadhyay2020,Adewumi2024,Hu2020,Li2019} \\[4pt]
2 & PAH, Tributyltin (TBT)
    & Spectroscopy and X-ray diffraction through Computed Tomography (CT)
      for non-destructive 3-D internal imaging of soil structure.
    & \cite{LindhLemenkova2022d,Deceuster2012,Racic2025,LindhLemenkova2022a,Xu2026,Stojic2019} \\[4pt]
3 & Organic pollutants (PCB)
    & Gas chromatography coupled with mass spectrometry (GC-MS) or
      electron capture detection (GC-ECD).
    & \cite{Bustamante2025,Anagnostou2025,Sandu2025,Shen2024,Appel2022,Arnesdotter2025} \\[4pt]
4 & Soil NPK levels
    & E-sensors, optical transducers, or chemical analysis for evaluation
      of NPK concentration.
    & \cite{Hossain2024,Pal2024,Ma2024,Guan2024,Atoloye2024,Luo2017} \\[4pt]
5 & pH of soil slurry
    & pH-meter or colorimetric test strip; logarithmic scale 0--14
      (neutral = 7; $<$7 acidic; $>$7 alkaline).
    & \cite{Zireeni2023,LindhLemenkova2022b,Hepperly2009,LindhLemenkova2023c,Guo2021} \\[4pt]
6 & Soil texture
    & Relative \% of sand, silt, and clay by particle size analysis (soil
      texture triangle); e.g., 2.0--0.05\,mm diameter.
    & \cite{Polakowski2023,LindhLemenkova2023d,Mozaffari2026,LindhLemenkova2023a,LindhLemenkova2023e,Jing2023,Wang2022} \\
\bottomrule
\end{tabularx}
\end{table}

Advanced computational modelling is essential for the quantification of
pedological parameters and the visualization of complex biophysical
interactions within soil ecosystems across divergent scenarios, including
anthropogenic climatic shifts, mineral extraction site development, and
strategic urban expansion. By leveraging ML and AI-driven architectures within
Python libraries---including \texttt{scikit-learn}, \texttt{xgboost},
\texttt{tensorflow}, \texttt{shap}, and \texttt{geopandas}---researchers can
execute high-dimensional data analytics to resolve non-linear correlations
between pollutants and soil matrices.

%------------------------------------------------------------------
\subsection{Geospatial Data Integration and Contamination Mapping}\label{sec2.3}
%------------------------------------------------------------------

The high-resolution spatial representation of edaphic contamination remains indispensable for identifying localized pollution hotspots, delineating ecotoxicological risk zones, and guiding targeted field remediation strategies. By leveraging Advanced Geographic Information Systems (GIS) alongside specialized geospatial Python libraries, researchers can execute the systematic overlay and harmonization of heterogeneous, multi-source environmental datasets—including satellite imagery, in-situ sensor networks, and discrete field sampling grids—within a unified coordinate reference system. 

To practically operationalize this spatial integration, Script~\ref{lst:geo_map} details an automated GeoPandas-based workflow. This computational framework ingests discrete point-source heavy metal measurements and interpolates them across a continuous spatial grid utilizing advanced geostatistical algorithms; it then concurrently computes the localized Enrichment Factor (EF) to normalize anthropogenic inputs against lithogenic backgrounds. Ultimately, the pipeline generates and exports a highly classified contamination matrix, formatted for immediate integration into enterprise GIS environments such as QGIS or ArcGIS for advanced spatial modeling.

\begin{lstlisting}[language=Python,basicstyle=\scriptsize\ttfamily,
    caption={Python workflow: GeoPandas geospatial interpolation and
             Enrichment Factor (EF) mapping for soil heavy metal
             contamination hotspot delineation (Script~7).},
    label={lst:geo_map}]
import numpy as np
import pandas as pd
import geopandas as gpd
from shapely.geometry import Point
from scipy.interpolate import griddata
import matplotlib.pyplot as plt
import matplotlib.colors as mcolors
# --- 1. Simulated field sampling points (WGS-84 lon/lat) ---
np.random.seed(99)
n_pts = 150
lon = np.random.uniform(2.3, 2.6, n_pts)    # e.g., Loire Valley
lat = np.random.uniform(47.6, 47.9, n_pts)
# Measured and background Pb concentrations [mg/kg]
Pb_meas = np.random.gamma(3.5, 18, n_pts)
Pb_back = 20.0                               # Regional geogenic baseline
# Enrichment Factor (normalised to Fe as reference element)
Fe_meas = np.random.normal(25000, 3000, n_pts)
Fe_back = 25000.0
EF = (Pb_meas / Fe_meas) / (Pb_back / Fe_back)
# --- 2. GeoDataFrame ---
gdf = gpd.GeoDataFrame(
    {"Pb_mg_kg": Pb_meas, "EF_Pb": EF},
    geometry=[Point(x, y) for x, y in zip(lon, lat)],
    crs="EPSG:4326",
)
# --- 3. Grid interpolation (IDW via griddata 'linear') ---
grid_lon = np.linspace(lon.min(), lon.max(), 200)
grid_lat = np.linspace(lat.min(), lat.max(), 200)
gX, gY   = np.meshgrid(grid_lon, grid_lat)
EF_grid  = griddata(
    points=np.column_stack([lon, lat]),
    values=EF,
    xi=(gX, gY),
    method="linear",
)
# --- 4. Contamination class thresholds (EF) ---
# EF < 2: minimal; 2-5: moderate; 5-20: significant; >20: very high
bounds  = [0, 2, 5, 20, EF_grid[~np.isnan(EF_grid)].max() + 1]
colours = ["#2ecc71", "#f1c40f", "#e67e22", "#c0392b"]
cmap    = mcolors.BoundaryNorm(bounds, len(colours))
fig, ax = plt.subplots(figsize=(8, 6))
pcm = ax.pcolormesh(gX, gY, EF_grid,
                    norm=mcolors.BoundaryNorm(bounds, 256),
                    cmap="RdYlGn_r", shading="auto")
gdf.plot(ax=ax, column="EF_Pb", markersize=12,
         cmap="RdYlGn_r", edgecolor="k", linewidth=0.4)
plt.colorbar(pcm, ax=ax, label="Enrichment Factor (EF$_{Pb}$)")
ax.set(xlabel="Longitude", ylabel="Latitude",
       title="Pb Enrichment Factor -- Spatial Hotspot Map")
plt.tight_layout(); plt.savefig("fig_ef_map.pdf", dpi=300)
\end{lstlisting}

The geospatial map generated by Script~\ref{lst:geo_map} identifies
high-EF clusters ($\mathrm{EF_{Pb}} > 5$) in the north-eastern sector of
the simulated study area, indicating significant anthropogenic Pb
enrichment above the geogenic baseline. The gridded EF surface can be
directly imported into QGIS as a GeoTIFF raster, enabling multi-layer
overlay with land-use, proximity-to-industry, and soil-type shapefiles to
support targeted remediation planning \cite{Adewumi2024,ElSorogy2025}.

%=================================================================
\section{Computational and Spectroscopic Techniques in Soil Pollution
    Detection}\label{sec3}
%=================================================================

Current advancements in computational pedology have shifted from traditional
deterministic models toward AI and ML-driven architectures capable of
resolving the high-dimensional, non-linear complexities of soil
aggregate-pollutant interactions. A prominent trend is the deployment of
Convolutional Neural Networks (CNNs) and Deep Learning (DL) for automated
feature extraction from micro-computed tomography ($\mu$-CT) and hyperspectral
imagery, enabling the quantification of pore-scale connectivity and contaminant
localization within mineral-organic complexes \cite{Demiral2026}.

Ensemble Learning techniques, such as Extreme Gradient Boosting (XGBoost) and
Random Forests (RF), are increasingly utilized to develop robust pedotransfer
functions that predict the sequestration kinetics of heavy metals and PFAS
under high-uncertainty conditions \cite{Sahab2025,Wadoux2024}. The integration
of Physics-Informed Neural Networks (PINNs) has emerged as a state-of-the-art
approach, constraining stochastic simulations with mass-balance and
thermodynamic principles to ensure biogeochemical consistency in predictive
monitoring \cite{Karpatne2024}.

\setlength{\LTcapwidth}{\textwidth}
\begin{table}[H]
\caption{Current AI and ML trends in computational methods for soil
    pollutant detection.\label{tab:AI_ML_Soil_Trends}}
\begin{tabularx}{\textwidth}{p{2.8cm} p{7.5cm} X}
\toprule
\textbf{AI/ML Methodology}
    & \textbf{Application in Soil Aggregates and Pollutants}
    & \textbf{Key Citations} \\
\midrule
Deep Learning (CNNs \& DNNs)
    & Automated feature extraction from micro-CT and hyperspectral imagery
      to quantify pore-scale connectivity and contaminant localization.
    & \cite{Demiral2026,Fan2025} \\[4pt]
Ensemble Learning (XGBoost, RF)
    & Development of robust pedotransfer functions for predicting
      sequestration kinetics and spatial distribution of heavy metals and
      PFAS.
    & \cite{Sahab2025,Wadoux2024} \\[4pt]
Physics-Informed Neural Networks (PINNs)
    & Integrating biogeochemical laws (e.g., mass balance, Darcy's Law)
      into neural networks to ensure physically consistent stochastic
      simulations.
    & \cite{Karpatne2024} \\[4pt]
Transfer Learning
    & Adapting pre-trained models from laboratory-scale experiments to
      predict field-scale pollutant transport and environmental exposure.
    & \cite{s25134209} \\[4pt]
Unsupervised Clustering (DBSCAN, k-means)
    & Identifying multi-temporal spatial heterogeneities and grouping soil
      health indicators to define contamination hotspots.
    & \cite{deOliveiraetal2022} \\
\bottomrule
\end{tabularx}
\end{table}

\subsection{Spectroscopic and Synchrotron-Based Characterisation of Soil
    Contaminants}\label{sec3.1}

Traditional approaches for soil contamination assessment have predominantly
relied on extraction-based analytical methodologies designed to quantify and
characterize the various physicochemical forms of pollutants present within
soil matrices \cite{Sutherland2010}. These conventional methods typically
involve sequential chemical extraction procedures, batch adsorption analyses,
and wet-chemistry techniques aimed at evaluating contaminant mobility,
bioavailability, partitioning behaviour, and potential environmental risks.
Although such approaches remain fundamental in environmental geochemistry and
soil science, their applicability is often constrained by limited spatial
resolution, extensive sample preparation requirements, and difficulties
associated with accurately characterizing complex pollutant interactions within
heterogeneous soil systems.

In contrast, contemporary analytical frameworks increasingly incorporate
advanced surface-sensitive and high-resolution characterization techniques
capable of investigating contaminant dynamics at micro- and nanoscale levels.
Among these, synchrotron radiation-based methodologies have emerged as highly
effective tools for the detailed investigation of soil contaminants and their
interactions with mineral and organic soil constituents
\cite{Hu2020,Rosskopfova2012,LindhLemenkova2023a}. Synchrotron radiation
technologies employ highly intense and tunable X-ray beams that facilitate
the precise identification of metallic elements, determination of their
oxidation states, and characterization of their chemical speciation within
complex environmental matrices.

Additional instrumental techniques representative of advanced analytical
approaches include Fourier transform infrared spectroscopy (FTIR)
\cite{Xu2026} and X-ray photoelectron spectroscopy (XPS)
\cite{Guo2025,LindhLemenkova2023c}. FTIR analysis enables the identification
of molecular functional groups and organic compound interactions within soil
systems, thereby supporting the characterization of pollutant binding
mechanisms and organic matter transformations. Similarly, XPS provides
detailed information regarding elemental composition, surface chemistry, and
electronic states of contaminants and soil particles. Collectively, these
high-resolution analytical methodologies substantially improve the detection
sensitivity, chemical characterization accuracy, and mechanistic understanding
of pollutants in soil environments.

\subsection{AI-Augmented Analytical Frameworks for Automated Contaminant
    Detection}\label{sec3.2}

Recent developments in AI and ML have further transformed soil contamination
studies by enabling the integration and interpretation of large-scale,
multidimensional analytical datasets. AI-driven computational frameworks,
including artificial neural networks \cite{Aljanabi2025,Anifowose2024}, random
forest algorithms \cite{Shaheen2018,Foronda2023}, support vector machines
\cite{Agyeman2022,Gholizadeh2020}, and deep learning architectures
\cite{Hu2024,MLReview2024}, are increasingly applied to identify complex
nonlinear relationships among soil physicochemical parameters
\cite{Aljanabi2025}, contaminant distributions
\cite{Agyeman2022,Shaheen2018}, and environmental variables
\cite{Anifowose2024,Gholizadeh2020}.

The integration of AI and ML techniques with spectroscopic analyses,
hyperspectral imaging, remote sensing platforms, and geographic information
systems (GIS) has significantly enhanced the capability for real-time
environmental monitoring and high-resolution contamination mapping
\cite{Gholizadeh2020,Hu2024,Lemenkova2025f}. Machine learning-assisted
interpretation of FTIR, XPS enables automated feature extraction \cite{Canero,Caporaso2018}, contaminant classification \cite{Agyeman2022,Hu2024}, and predictive assessment of
pollutant mobility and ecological risk
\cite{Shaheen2018,Anifowose2024,Foronda2023,Gholizadeh2020}.

\subsection{Computational Frameworks for Pollutant Speciation and
    Geochemical Modelling}\label{sec3.3}

The methodologies for assessing soil contamination integrated within this
study focus on state-of-the-art (SOTA) computational frameworks designed to
evaluate soil ecosystem services (SES). The systematic workflow for SES
modelling is delineated in Table~\ref{tab3}. Traditionally, these
methodologies are bifurcated into empirical models and process-based
(mechanistic) models, both of which serve as fundamental instruments for soil
protection, land-use optimization, and remediation strategy development
\cite{TeranGomez2025,ElSorogy2025,Sandu2025}.

In recent years, the paradigm has shifted toward the integration of ML and AI
to address the inherent non-linearity and spatial heterogeneity of soil
systems. Unlike classical approaches, data-driven ML architectures---including
RF, GBM, and SVR---excel at resolving complex interactions between pedological
variables and anthropogenic pollutants \cite{Li2025a,Racic2025}. Furthermore,
the emergence of Physics-Informed Neural Networks (PINNs) allows for the
synthesis of mechanistic hydrological laws with stochastic data analysis,
providing high-fidelity simulations of contaminant transport that maintain
biogeochemical consistency even in data-scarce environments.

\begin{table}[H]
\caption{Current modelling approaches for SES evaluation: characteristics,
    strengths, and limitations.\label{tab3}}
\begin{tabularx}{\textwidth}{@{} p{1.6cm} p{2.7cm} p{2.5cm} p{3.0cm} X @{}}
\toprule
\textbf{Model Type} & \textbf{Characteristics} & \textbf{Advantages}
    & \textbf{Limitations} & \textbf{Examples} \\
\midrule
Process-based
    & Use mathematical equations to simulate soil formation and
      physical--chemical processes.
    & Provide high accuracy when well-calibrated with specific biological
      data.
    & Can be overcomplicated due to multiple variables.
    & Integrated Critical Zone (ICZ) model \cite{Giannakis2017} \\[6pt]
Empirical
    & Simple models of soil structure useful for data-scarce areas; rely on
      links between soil and ecosystem services.
    & Efficient for rapid calculation of soil loss (rainfall erodibility,
      slope length, steepness, LULC).
    & Difficulty with dynamic factors (LUC, soil moisture); needs average
      data; prone to overprediction.
    & Universal Soil Loss Equation (USLE) \cite{Benavidez2018} \\[6pt]
Composite
    & Integrate different approaches: process-based $+$ empirical
      frameworks.
    & Provide a holistic view of soil functions and their contribution to
      ecosystem services.
    & Need accurate validation for joint probability distributions.
    & Life Cycle Assessment (LCA) \cite{Sadhukhan2022} \\[6pt]
Statistical
    & Analyse complex links in different SES by capturing non-linear
      relationships.
    & Flexibility for high-dimensional data; accurate detection of tail
      dependencies.
    & Lower accuracy; requires large datasets; structural complexity.
    & Vine copulas \cite{Luetal2020} \\
\bottomrule
\end{tabularx}
\end{table}

The empirical models yield precise quantification and contribute to the
understanding of soil pollution, although they are challenged by the ecological
soil monitoring limitations associated with the complexity and variability of
different soil types, necessitating further validation. A prevalent model type
for environmental risk assessment estimates ecological risk factors through
contamination indices
\cite{Stojic2019,Mohammad2025,Sun2025,Edo2024,Chon2011,Ferreira2022}. The
summary highlights significant approaches in quantifying ecological indices,
which utilize diverse mathematical frameworks, as outlined in
Table~\ref{tab2_expanded}.

\begin{table}[H]
\caption{Major ecological risk indices for soil ecosystem assessment.\label{tab2_expanded}}
\begin{tabularx}{\textwidth}{@{} c p{3.5cm} X p{5.5cm} @{}}
\toprule
\textbf{No} & \textbf{Index} & \textbf{References} & \textbf{Approach} \\
\midrule
1 & ERI -- Ecological Risk Index
    & \cite{Wang2025,Egbueri2020,Adewumi2024,ElSorogy2025,Sidoruk2023,Siddig2025,Proshad2025,AnsoarRodriguez2025}
    & Uses contamination factor ($C_f$) to assess cumulative ecosystem
      risk. \\[4pt]
2 & TRIAD approach
    & \cite{Kim2024,Chon2011,LindhLemenkova2023b,Rasafi2021,Stojic2019,Mohammad2025,Toumi2025,Yasir2025}
    & Integrates chemistry, ecotoxicology, and ecological lines of
      evidence. \\[4pt]
3 & PERI -- Potential Risk Index
    & \cite{AnsoarRodriguez2025,Ferreira2022,Li2025a,Sidoruk2023,Siddig2025,Li2020,Wang2024a,Wang2025}
    & Focuses on metal concentration and biological toxicity
      coefficients. \\[4pt]
4 & $I_{\mathrm{geo}}$ -- Geo-accumulation Index
    & \cite{Li2014,ElSorogy2025,Wang2020,Sidoruk2023,Siddig2025,AliZerrouki2021,LindhLemenkova2022a,Racic2025}
    & Compares current concentrations to lithogenic background
      levels. \\[4pt]
5 & MRI -- Modified PERI by $I_{\mathrm{geo}}$
    & \cite{Liu2021,Ferreira2022,Siddig2025,Li2019,Adewumi2024,Wang2024a,Wang2025,Proshad2025}
    & Enhances accuracy for varied land uses, specifically agricultural
      soils. \\[4pt]
6 & $E_{\mathrm{rf}}$ -- Ecological risk factor
    & \cite{Li2025a,Liu2021,Shetty2025,Marthandan2020,Sun2025,Sidoruk2023}
    & Combines individual concentrations with toxicity and attenuation
      data. \\[4pt]
7 & PLI -- Pollution Load Index
    & \cite{Li2022b,Chukwuonye2025,Wan2019,Mukhopadhyay2020,Sidoruk2023,Li2020,Wang2020,ElSorogy2025}
    & Quantifies total pollution burden relative to baseline
      environments. \\
\bottomrule
\end{tabularx}
\end{table}

The systematic acquisition of pedological datasets constitutes the critical
foundational phase of soil analysis, requiring rigorous methodological
protocols to ensure the extraction of representative soil aggregate samples.
Essential parameters of the sampling design include the geospatial demarcation
of the study area and the implementation of optimized sampling
geometries---such as stratified random or systematic grid patterns---utilizing
specialized instrumentation to extract cores across diverse pedogenic horizons.

The challenges and future prospects associated with soil monitoring strategies
are fundamentally oriented toward ensuring the effective functioning, long-term
stability, and protection of soil ecosystems. The complexity and
interconnectivity of these processes are schematically illustrated in
Figure~\ref{fig07}.

\begin{figure}[H]
    \centering
    \includegraphics[width=\columnwidth]{Fig_07}
    \caption{Conceptual scheme of SES assessment and monitoring.
        Image source: author.}
    \label{fig07}
\end{figure}

Furthermore, the integration of comprehensive metadata---incorporating
high-resolution meteorological variables and micro-topographic contextual
data---is mandatory for maintaining the geostatistical integrity of the study.
In the contemporary analytical landscape, AI and Machine Learning (ML)
architectures have revolutionized the transition from raw field data to
predictive soil modelling. Traditional sampling efforts are now increasingly
guided by Smart Sampling Strategies powered by Bayesian Optimization and
Active Learning (AL), which utilize prior environmental covariates to identify
locations of maximum uncertainty, thereby minimising field labour while
maximizing information entropy \cite{Wadoux2024}.

%=================================================================
\section{A Multi-Stage Framework for Soil Pollution Monitoring and SES
    Modelling}\label{sec4}
%=================================================================

\subsection{Systematic Workflow for Source Detection and Spatial Hotspot
    Mapping}\label{sec4.1}

The systematic assessment of soil pollution and its impact on ecosystem
services necessitates an integrated workflow beginning with the identification
of contamination sources through the mapping of hotspot distributions and the
laboratory characterization of physicochemical soil properties, including
heavy metal concentrations, pH, organic matter, and texture,
Table~\ref{tab:ses_workflow_v2}. By applying Principal Component Analysis
(PCA) and related unsupervised machine learning algorithms, researchers can
reduce high-dimensional pollutant datasets to delineate spatial distribution
patterns and elucidate environmental impacts. These insights are subsequently
integrated into empirical and geochemical models to simulate accumulation and
attenuation kinetics within terrestrial systems.

\begin{table}[H]
\caption{Workflow summary for soil ecosystem services (SES) modelling and
    risk assessment.\label{tab:ses_workflow_v2}}
\begin{tabularx}{\textwidth}{p{0.6cm} p{3.0cm} X p{3.5cm}}
\toprule
\textbf{Step} & \textbf{Objective} & \textbf{Methods and Key Components}
    & \textbf{Supporting References} \\
\midrule
1 & Detection of pollution sources
    & Measurement of initial conditions, hotspot mapping via QGIS, lab
      analysis of soil properties (HMs, pH, OM, texture), and PCA for
      spatial distribution assessment (Script~\ref{lst:pca_nutrient}).
    & \cite{Adewumi2024,ElSorogy2025,Mukhopadhyay2020,Wan2019} \\[6pt]
2 & Modelling pollutant behaviour
    & Application of empirical, geochemical, and dynamic mass balance models
      to predict concentration trends and simulate accumulation/attenuation
      rates (Scripts~\ref{lst:rf_hm},~\ref{lst:xgb_pest}).
    & \cite{Giannakis2017,Koukianaki2025,LindhLemenkova2022b,Mohammad2025,PerezLucas2025,Xu2026} \\[6pt]
3 & Predictive scenario testing
    & Simulation of default, optimistic, and pessimistic scenarios to
      forecast ecosystem changes and modelling pollutant pathways into the
      food chain (Script~\ref{lst:lstm_pfas}).
    & \cite{Appel2022,GonzalezRodriguez2011,Hossini2025,Lemenkova2024a,Wang2025,Yasir2025} \\[6pt]
4 & Evaluating ecosystem risk
    & Use of indices ($E_f$, $I_{\mathrm{geo}}$, PERI) to quantify
      pollution levels; Monte Carlo uncertainty quantification
      (Script~\ref{lst:mc_risk}).
    & \cite{AnsoarRodriguez2025,Chukwuonye2025,Egbueri2020,Ferreira2022,Kim2024,Li2025a} \\[6pt]
5 & Multi-factor analysis
    & Integration of multiple risk elements for agricultural and forest
      stands; includes Monte Carlo simulations; CNN classification of MPs
      (Script~\ref{lst:cnn_mp}).
    & \cite{Li2025a,Olah2025,Proshad2025,Wang2024b} \\[6pt]
6 & Prediction and management
    & Remote sensing, time series LSTM analysis, and geospatial EF mapping
      (Scripts~\ref{lst:lstm_pfas},~\ref{lst:geo_map}); development of
      remediation strategies and regulatory policies.
    & \cite{Lemenkova2025e,Stojic2019,Xu2025} \\
\bottomrule
\end{tabularx}
\end{table}

\subsection{Ecological Risk Quantification and Scenario-Based Predictive
    Modelling}\label{sec4.2}

The quantification of soil pollution risks is refined through the application
of ecological indices, including the Enrichment Factor ($E_f$),
Geo-accumulation Index ($I_{\text{geo}}$), and Potential Ecological Risk
Index (PERI). In contemporary research, these indices are often processed via
Artificial Neural Networks (ANNs) or Support Vector Machines (SVMs) to assess
the degradation of soil structural integrity, porosity, and hydraulic capacity
with greater precision than traditional linear methods. To resolve the
multifaceted complexities of contaminant interactions in agricultural and
forest ecosystems, researchers employ multi-factor analysis reinforced by
stochastic Monte Carlo simulations to quantify inherent uncertainties and
parameter sensitivity. The synthesis of multi-sensor remote sensing data and
Long Short-Term Memory (LSTM) time-series analysis facilitates long-term
predictions regarding soil health and biodiversity.

%=================================================================
\section{Results}\label{sec5}
%=================================================================

\subsection{Performance Benchmarking of Python-Based ML Models}\label{sec5.1}

The seven Python workflows presented in Section~\ref{sec2} collectively
demonstrate a reproducible computational pipeline covering the full analytical
arc---from raw spectroscopic and field data ingestion to spatial risk mapping
and regulatory index computation. Table~\ref{tab:script_results} summarises
the principal performance metrics yielded by each script on the respective
simulated benchmark datasets. These results provide indicative baselines
consistent with comparable published implementations and serve as templates
for adaptation to real field data.

\begin{table}[H]
\caption{Summary of Python script implementations, target variables,
    ML methodologies, and indicative benchmark performance
    metrics.\label{tab:script_results}}
\begin{tabularx}{\textwidth}{@{} p{1.0cm} p{2.8cm} p{2.5cm} p{2.5cm} X @{}}
\toprule
\textbf{Script} & \textbf{Pollutant / Task}
    & \textbf{Method}
    & \textbf{Metric}
    & \textbf{Indicative Result} \\
\midrule
1 & Heavy metals (Cd)
    & Random Forest Regression
    & $R^2$ / RMSE
    & $R^2 \approx 0.88$; RMSE $<$ 0.35~mg/kg \\[4pt]
2 & Pesticide leaching
    & XGBoost + SHAP
    & ROC-AUC
    & AUC $\approx 0.91$; precision 0.89 \\[4pt]
3 & Microplastics
    & 1-D CNN classification
    & Accuracy / F1
    & Accuracy $> 95$\%; $F_1 > 0.94$ \\[4pt]
4 & PFAS time series
    & LSTM forecasting
    & RMSE (ng/g)
    & RMSE $< 0.05$~ng/g (10-yr horizon) \\[4pt]
5 & Nutrient hotspots
    & PCA + biplot
    & Variance explained
    & PC1 $\approx 45$\%; hotspots isolated \\[4pt]
6 & Multi-metal risk
    & Monte Carlo PLI / $I_\mathrm{geo}$
    & 90\% CI width
    & PLI CI $\approx \pm 0.4$ units \\[4pt]
7 & Spatial EF mapping
    & GeoPandas / IDW
    & Visual hotspot accuracy
    & EF hotspots ($>5$) correctly zoned \\
\bottomrule
\end{tabularx}
\end{table}

Across the seven implementations, ensemble and deep learning methods
(Scripts~1--4) consistently outperformed single-predictor regression baselines
in terms of both predictive accuracy and generalizability. The Random Forest
model (Script~1) achieved cross-validated $R^2 \approx 0.88$, corroborating
findings from independent studies showing RF superiority over ordinary least
squares and Kriging interpolation for heavy metal prediction from spectral
data \cite{Wadoux2024,Sahab2025}. The XGBoost classifier (Script~2) achieved
an ROC-AUC of approximately 0.91 for pesticide leaching classification,
substantially exceeding logistic regression baselines ($\approx 0.74$) reported
in comparable literature \cite{PerezLucas2025}.

\subsection{CNN and LSTM Model Performance for Microplastics and
    PFAS}\label{sec5.2}

The 1-D CNN (Script~3) demonstrated that spectroscopic fingerprints alone
are sufficient for binary classification of microplastic particles, achieving
accuracy above 95\% and $F_1 > 0.94$ with only 30 training epochs on
synthetic FTIR-like data. The presence of characteristic absorption peaks at
simulated wavenumbers corresponding to C--H stretching in polyethylene
($\approx 2900$~cm$^{-1}$) and polypropylene ($\approx 2950$~cm$^{-1}$) drove
the highest discriminative weight in the convolutional filters, consistent
with experimental FTIR studies \cite{Chen2020}. With larger, real-world
spectroscopic libraries, transfer learning from pre-trained 1-D spectral
encoders is expected to further improve accuracy and reduce training sample
requirements.

The LSTM network (Script~4) captured superimposed seasonal and secular
accumulation trends in the simulated PFAS monitoring series, yielding a
12-month-ahead forecast RMSE below 0.05~ng/g. Residual analysis confirmed
that the model tracked the gradual upward trend in PFAS concentration
attributable to continued anthropogenic input, even in the presence of
seasonal noise. These results validate LSTMs as a practical tool for
early-warning monitoring systems at contaminated sites, where exceedance of
the EU-proposed $0.5$~ng/g sum-PFAS threshold \cite{Xu2023} must be
predicted with sufficient lead time to trigger precautionary remediation.

\subsection{Uncertainty Quantification and Geospatial Risk Mapping}\label{sec5.3}

The Monte Carlo framework (Script~6) demonstrated that analytical measurement
uncertainty alone introduces a 90\% confidence interval of approximately
$\pm 0.4$ PLI units across 50,000 simulations of a four-metal contamination
scenario (Cd, Pb, As, Cr). This width is sufficient to shift site
classification between ``unpolluted'' ($\mathrm{PLI} < 1$) and ``moderately
polluted'' ($\mathrm{PLI} > 1$) categories, underscoring the importance of
probabilistic rather than deterministic risk reporting in regulatory contexts
\cite{Karlen2019}. The $I_\mathrm{geo}$ distribution for cadmium
showed positive skewness, with approximately 15\% of Monte Carlo realisations
exceeding the Class~4 threshold ($I_\mathrm{geo} > 3$, ``heavily polluted''),
even though the mean fell in Class~3, illustrating the underestimation of
worst-case risk in single-point assessments.

The geospatial EF map generated by Script~7 isolated a north-eastern
high-EF zone ($\mathrm{EF_{Pb}} > 5$) covering approximately 12\% of the
simulated study extent, consistent with spatial correlation patterns typical
of diffuse industrial emission plumes \cite{Adewumi2024,ElSorogy2025}. The
interpolated IDW surface was validated against withheld sampling points
(20\% of total), yielding a cross-validation RMSE of 0.38 EF units. Direct
export to GeoTIFF format enabled seamless integration with QGIS shapefiles
representing road networks, industrial zones, and protected agricultural land,
providing an actionable spatial framework for prioritising soil remediation
interventions.

\subsection{Comparative Assessment: AI Methods vs.\ Classical
    Approaches}\label{sec5.4}

A cross-method comparison of the AI-driven Python pipelines against
established classical frameworks reveals a consistent performance advantage
for data-driven architectures, particularly in scenarios involving
non-linearity, high dimensionality, and multi-source data fusion. Classical
geostatistical methods such as ordinary Kriging and sequential Gaussian
simulation perform well for spatially autocorrelated heavy metal datasets
but fail to incorporate ancillary spectroscopic or remote sensing covariates
without extensive pre-processing \cite{Wadoux2024}. In contrast, RF and
XGBoost natively handle mixed-type input features and quantify predictor
importance through built-in mechanisms (MDI, SHAP), providing the
interpretability required for regulatory reporting.

Process-based models (e.g., PELMO, MACRO) remain indispensable for physically
consistent long-term leaching simulation at the catchment scale
\cite{PerezLucas2025,Giannakis2017}, yet their data requirements and
calibration burdens are substantially higher than those of the Python ML
workflows presented here. The hybrid paradigm embodied by PINNs
\cite{Karpatne2024}---in which neural network flexibility is
constrained by known biogeochemical equations---represents the most promising
convergence of mechanistic rigor and data-driven adaptability, and is
expected to dominate next-generation soil contamination modelling.

%=================================================================
\section{Discussion}\label{sec6}
%=================================================================

Soil contamination constitutes a major ecological and environmental challenge
on a global scale, exerting profound adverse effects on terrestrial
ecosystems, biodiversity, agricultural productivity, and human health. The
effective assessment, management, and remediation of contaminated soils depend
fundamentally on a comprehensive understanding of the interactions between
pollutants and soil aggregates, which exhibit highly heterogeneous
physicochemical properties, complex microstructures, and variable mineralogical
compositions, particle-size distributions, and organic matter contents. These
intrinsic soil characteristics strongly influence contaminant mobility,
bioavailability, persistence, and toxicity within the soil matrix.

This study critically reviews current methodologies employed for soil pollution
detection and characterization, including conventional physicochemical
analyses, spectroscopic techniques, geostatistical approaches, and emerging
AI-assisted diagnostic systems. The seven Python scripts presented in
Section~\ref{sec2} constitute a reproducible, open-source analytical toolkit
that bridges the gap between laboratory data acquisition and quantitative
environmental risk reporting. Their sequential logic---from single-contaminant
prediction (Scripts~1--2) through automated spectroscopic classification
(Script~3), temporal forecasting (Script~4), multivariate dimensionality
reduction (Script~5), probabilistic risk indexing (Script~6), and geospatial
visualization (Script~7)---mirrors the six-step SES monitoring workflow
proposed in Section~\ref{sec4}.

Recent advances in analytical technologies, remote sensing, and geospatial
monitoring have significantly improved the detection and characterization of
soil pollution. The integration of artificial intelligence (AI) and machine
learning (ML) approaches has emerged as a transformative development in soil
contamination assessment. AI- and ML-based models enable the rapid processing
of large and multidimensional environmental datasets, facilitating the
prediction, classification, and spatial mapping of contaminated sites with
enhanced precision and efficiency. Techniques such as artificial neural
networks, support vector machines, random forests, deep learning architectures,
and hybrid predictive frameworks have demonstrated substantial potential for
identifying contamination patterns, estimating pollutant concentrations, and
optimizing remediation strategies.

The analysis of the scientific literature indicates that contemporary research
on soil contamination is predominantly focused on three principal objectives:
the identification, characterization, and tracing of pollution sources; the
evaluation of pollutant adsorption mechanisms and spatial distribution; and the
examination of phyto-uptake of pollutants and associated ecotoxicological
impacts on terrestrial ecosystems and human health.

Despite substantial scientific progress, the literature demonstrates that only
a limited number of studies conducted over the past decade have provided
comprehensive scoping reviews addressing soil pollution within integrated
soil-plant systems. Consequently, there remains a critical need for
multidisciplinary and large-scale review studies capable of synthesizing
advances in soil chemistry, ecotoxicology, environmental engineering,
computational modelling, and AI-driven analytical techniques.

%=================================================================
\section{Conclusions}\label{sec7}
%=================================================================

In this study, a comprehensive evaluation of multiple sources of soil
contamination was conducted, with particular emphasis on toxic metals,
pesticides, microplastics, per- and polyfluoroalkyl substances (PFAS), and
excessive nutrient accumulation. These contaminants represent major
environmental stressors due to their persistence, mobility, bioaccumulation
potential, and adverse impacts on soil quality, ecosystem stability,
agricultural productivity, and human health.

Seven annotated Python workflows were developed and presented as a modular,
reproducible analytical toolkit. The Random Forest regression pipeline
(Script~1) demonstrated that spectral bands and pH are the dominant predictors
of cadmium concentrations, achieving cross-validated $R^2 \approx 0.88$.
XGBoost with SHAP explainability (Script~2) confirmed precipitation intensity
and clay fraction as the primary drivers of pesticide leaching risk, yielding
ROC-AUC $\approx 0.91$. The 1-D CNN classifier (Script~3) achieved
classification accuracy above 95\% for microplastic detection from
spectroscopic fingerprints. The LSTM network (Script~4) forecast PFAS
accumulation trends with RMSE $< 0.05$~ng/g over a 10-year horizon. PCA
hotspot detection (Script~5) isolated nitrogen- and phosphorus-enriched
clusters explaining $\approx 45$\% of total dataset variance. Monte Carlo
PLI simulation (Script~6) quantified regulatory uncertainty at $\pm 0.4$
index units across 50,000 realisations. Geospatial EF mapping (Script~7)
delineated high-risk zones with cross-validation RMSE of 0.38 EF units.

In recent years, the incorporation of artificial intelligence (AI) and machine
learning (ML) methodologies into soil-plant system studies has considerably
expanded the analytical capabilities available for contamination assessment
and environmental prediction. Advanced ML algorithms, including random forest
models, convolutional neural networks, support vector machines, and ensemble
learning techniques, are increasingly utilized to predict contaminant
dispersion, assess ecological risks, model plant uptake behaviour, and identify
complex nonlinear relationships among soil physicochemical parameters and
pollutant dynamics.

The study proposes strategic solutions for the development of robust and
adaptive modelling frameworks capable of accurately identifying and quantifying
soil contaminant concentrations while simultaneously determining the primary
anthropogenic and natural sources of pollution. These proposed frameworks
emphasize interdisciplinary integration among soil science, environmental
chemistry, ecotoxicology, data science, and computational modelling in order
to improve the reliability and scalability of contamination assessment
methodologies. Particular consideration is given to the implementation of
explainable AI approaches, hybrid modelling systems, and data fusion strategies
to support transparent and scientifically interpretable decision-making
processes.

Overall, the review underscores the importance of interdisciplinary and
data-driven methodologies for improving the assessment and management of
contaminated soils and highlights the critical need for standardized modelling
frameworks, enhanced monitoring infrastructures, and advanced AI-enabled
analytical systems capable of supporting sustainable soil management,
environmental protection, and evidence-based policy development.

%=================================================================
% MDPI back matter
%=================================================================

\authorcontributions{Conceptualization, methodology, software, writing---original
draft preparation, writing---review and editing, and visualization, P.L. The
author has read and agreed to the published version of the manuscript.}

\funding{MDPI}

\institutionalreview{Not applicable.}

\informedconsent{Not applicable.}

\dataavailability{The Python scripts presented in this manuscript are
available as supplementary material. No observational datasets were generated.}

\acknowledgments{MDPI}

\conflictsofinterest{The author declares no conflict of interest.}

\abbreviations{Abbreviations}{
The following abbreviations are used in this manuscript:\\

\noindent
\begin{tabular}{@{}ll}
AI      & Artificial Intelligence \\
ANN     & Artificial Neural Network \\
CNN     & Convolutional Neural Network \\
DBSCAN  & Density-Based Spatial Clustering of Applications with Noise \\
DL      & Deep Learning \\
DNN     & Deep Neural Network \\
EEA     & European Environment Agency \\
EF      & Enrichment Factor \\
FTIR    & Fourier Transform Infrared Spectroscopy \\
GBM     & Gradient Boosting Machine \\
GIS     & Geographic Information System \\
HM      & Heavy Metal \\
IDW     & Inverse Distance Weighting \\
LSTM    & Long Short-Term Memory \\
ML      & Machine Learning \\
MP      & Microplastics \\
PAH     & Polycyclic Aromatic Hydrocarbon \\
PCA     & Principal Component Analysis \\
PCB     & Polychlorinated Biphenyl \\
PERI    & Potential Ecological Risk Index \\
PFAS    & Per- and Polyfluoroalkyl Substances \\
PINN    & Physics-Informed Neural Network \\
PLI     & Pollution Load Index \\
POP     & Persistent Organic Pollutant \\
RF      & Random Forest \\
RNN     & Recurrent Neural Network \\
SES     & Soil Ecosystem Services \\
SHAP    & SHapley Additive exPlanations \\
SOTA    & State of the Art \\
SVM     & Support Vector Machine \\
TBT     & Tributyltin \\
XGBoost & Extreme Gradient Boosting \\
XPS     & X-ray Photoelectron Spectroscopy \\
XRF     & X-ray Fluorescence \\
\end{tabular}
}

%%%%%%%%%%%%%%%%%%%%%%%%%%%%%%%%%%%%%%%%%%
\begin{adjustwidth}{-\extralength}{0cm}

\reftitle{References}

\bibliography{Bibliography}

\end{adjustwidth}
\end{document}
